\documentclass{article}
\usepackage[paperwidth=160mm,paperheight=154mm,margin=0pt]{geometry}
\usepackage[UTF8,fontset=windows]{ctex}
\usepackage{tikz}
\usepackage{graphicx}
\usetikzlibrary{arrows.meta,positioning,calc}
\begin{document}
\pagestyle{empty}\noindent
\begin{tikzpicture}[x=1mm,y=1mm, every node/.style={font=\sffamily}]
\definecolor{ink}{HTML}{26374F}
\definecolor{cyan}{HTML}{E3F3F5}
\definecolor{pink}{HTML}{F9E8EB}
\definecolor{green}{HTML}{E7F4EB}
\definecolor{cream}{HTML}{FFF3DD}
\definecolor{blue}{HTML}{4A9DAD}
\fill[white] (0,0) rectangle (160,154);
\fill[cyan] (0,113) rectangle (160,154);
\fill[pink] (0,71) rectangle (160,110);
\fill[green] (0,0) rectangle (160,68);
\node[anchor=west,text=ink,font=\sffamily\bfseries\fontsize{12}{14}\selectfont] at (5,148) {问题三  共享只读L2与成对反事实};
\node[anchor=west,text=ink,font=\sffamily\bfseries\fontsize{10}{12}\selectfont] at (5,137) {01 对象与机理};
\node[anchor=west,text=ink,font=\sffamily\fontsize{9}{11}\selectfont] at (6,125) {COPY\_IN按逻辑张量编号查询共享L2};

\node[draw=blue!60,fill=white,rounded corners=1mm,minimum width=25mm,minimum height=12mm,font=\fontsize{8}{9}\selectfont] (query) at (83,130) {COPY\_IN查询};
\node[draw=blue!60,fill=cream,rounded corners=1mm,minimum width=22mm,minimum height=12mm,font=\fontsize{8}{9}\selectfont] (cache) at (122,137) {L2读池};
\node[draw=blue!60,fill=white,rounded corners=1mm,minimum width=22mm,minimum height=12mm,font=\fontsize{8}{9}\selectfont] (ddr) at (122,121) {DDR读池};
\draw[blue,thick,-{Latex[length=2mm]}] (query.north east)--node[above,font=\fontsize{7}{8}\selectfont]{命中}(cache.west);
\draw[orange!85!black,thick,-{Latex[length=2mm]}] (query.south east)--node[below,font=\fontsize{7}{8}\selectfont]{未命中}(ddr.west);

\node[anchor=west,text=ink,font=\sffamily\bfseries\fontsize{10}{12}\selectfont] at (5,103) {02 方法与约束};
\node[rounded corners=1mm,draw=blue!50,fill=white,minimum width=43mm,minimum height=19mm,align=center,text=ink,font=\sffamily\fontsize{8.6}{10.5}\selectfont] (a) at (26,89) {命中：Cache读池\\FIFO次序不刷新};
\node[rounded corners=1mm,draw=blue!50,fill=white,minimum width=43mm,minimum height=19mm,align=center,text=ink,font=\sffamily\fontsize{8.6}{10.5}\selectfont] (b) at (80,89) {未命中：DDR读池\\完成后按容量写入};
\node[rounded corners=1mm,draw=blue!50,fill=white,minimum width=43mm,minimum height=19mm,align=center,text=ink,font=\sffamily\fontsize{8.6}{10.5}\selectfont] (c) at (134,89) {同方案场景B对照\\隔离Cache机制收益};
\draw[blue,thick,-{Latex[length=2mm]}] (a.east) -- (b.west);
\draw[blue,thick,-{Latex[length=2mm]}] (b.east) -- (c.west);
\node[anchor=west,text=ink,font=\sffamily\fontsize{8.2}{10}\selectfont] at (5,74) {$R_k=T_k^{(0)}/T_k^{(L2)}$；命中字节率按官方口径};
\node[anchor=west,text=ink,font=\sffamily\bfseries\fontsize{10}{12}\selectfont] at (5,62) {03 实测结果与核验};
\node[anchor=west,text=ink,font=\sffamily\bfseries\fontsize{8.8}{10.5}\selectfont] at (5,52) {100例五核同方案收益均值 $1.0260$；58例获益};
\node[anchor=west,text=ink,font=\sffamily\fontsize{8.3}{10}\selectfont] at (5,44) {正式默认配置复核；六例54格非官方敏感性单独标识};
\node[anchor=west] at (5,21) {\includegraphics[width=72mm,height=30mm,keepaspectratio]{E:/20hw_myskill_distribution/workshop_v1/求解/正式图件/q3-F02.pdf}};
\node[anchor=east] at (155,21) {\includegraphics[width=72mm,height=30mm,keepaspectratio]{E:/20hw_myskill_distribution/workshop_v1/求解/正式图件/q3-F03.pdf}};
\node[anchor=west,text=ink,font=\sffamily\fontsize{7.5}{9}\selectfont] at (5,4) {细节参见正文图 q3-F02、q3-F03；缩图仅作证据导航。};
\end{tikzpicture}
\end{document}
